\er{old abstract, to be revised.}
Broadcast is a fundamental operation in Mobile Ad-Hoc Networks (MANETs).
A large variety of broadcast protocols have been proposed to deal with network dynamics, maximize coverage and minimize the impact of collisions.
These protocols follow either an unstructured approach based on the control of a flooding process or a more structured approach involving the creation of a dissemination overlay.
The selection of the appropriate broadcast protocol for a MANET is a difficult task.
Each protocol is better adapted to particular configurations of nodes, depending in particular on their density over space or their mobility.
For heterogeneous node configurations, a single protocol is only adequate for some regions of the network offering poor performance elsewhere.
This paper presents a novel adaptive approach to broadcast in heterogeneous MANETs.
This approach allows MANET nodes to pick dynamically the most appropriate broadcast protocols for different regions of the network.
We perform simulations and evaluate our approach using a mobility model that mimics how humans move to zones of interest, where overlays of nodes emerge as a result of applying our adaptive solution.
%
Our evaluation results show that our adaptive approach leads to better performance and lower network consumption.
